\section*{Chapitre \ref{ch:rs:portfolio-construction-ii} --- Construction de portefeuille II}

\begin{solution}{exo:rs:portfolio-construction-ii:1}
$\sqrt{0.5^2 \times 0.04 + 0.5^2 \times 0.01} = 11.2\%$. Sans corrélation, les \omterm{def:rs:portfolio-construction-ii:erc}{contributions au risque} égales valent $w_i\sigma_i^2$, égales après normalisation, donc
$w_i \propto 1/\sigma_i$~: $1/3$ et $2/3$.
\end{solution}

\begin{solution}{exo:rs:portfolio-construction-ii:2}
Sans vues, la moyenne a posteriori vaut $\Pi = \gamma\Sigma w_m$, et l'optimum moyenne--variance sur celle-ci (de covariance $\Sigma$) vaut $(\gamma\Sigma)^{-1}\gamma\Sigma w_m = w_m$~;
avec la covariance a posteriori $(1+\tau)\Sigma$, il vaut $w_m/(1+\tau)$, soit le marché remis à l'échelle.
\end{solution}

\begin{solution}{exo:rs:portfolio-construction-ii:3}
$1/(1 + \phi \times 0.78)$~: 0.56 pour le retournement, 0.987 pour la dérive post-annonce, 0.999 pour la dérive persistante.
\end{solution}

\begin{solution}{exo:rs:portfolio-construction-ii:4}
La réponse moyenne est fixée par l'inverse de la covariance et par $\gamma$, que le terme robuste modifie peu. Les réponses les plus fortes viennent des titres que l'optimiseur
juge quasi interchangeables, où une faible variation d'alpha reporte du poids de l'un à l'autre~; la pénalité sur $\sqrt{w^\top\Omega w}$ rend coûteux de se concentrer sur l'un ou
l'autre, si bien que le portefeuille détient les deux et que le basculement est amorti (de 0.33~\% à 0.12~\% au plus).
\end{solution}

\begin{solution}{exo:rs:portfolio-construction-ii:5}
Il minimise la variance sous une contrainte sur la diversité des poids ($\sum\ln w_i$ au-dessus d'une borne)~: sans contrainte, la solution est la variance minimale,
avec une contrainte infiniment serrée il donne des poids égaux, et le \omterm{def:rs:portfolio-construction-ii:erc}{portefeuille à contributions au risque égales} est un point de ce chemin (Maillard, Roncalli et Teïletche).
\end{solution}

\begin{solution}{exo:rs:portfolio-construction-ii:6}
Les deux échangent chaque jour 56~\% de l'écart, mais l'ajustement partiel vise le portefeuille de Markowitz, dominé par le retournement qui s'éteint en un jour~: il continue de payer pour se rapprocher
d'une cible qui aura bougé entre-temps. Le \omterm{def:rs:portfolio-construction-ii:aim}{portefeuille cible} (aim portfolio) pondère le retournement par 0.56 et les signaux lents presque par un, si bien que la cible elle-même bouge moins.
\end{solution}

\begin{solution}{exo:rs:portfolio-construction-ii:7}
Pour $\lambda = 1\,000$~: Markowitz quotidien $-22.62$, cible 3.15, partiel 1.08, mensuel 2.55. Pour $\lambda = 10\,000$~: la règle de cible 1.61, le rééquilibrage mensuel $-0.68$, les
autres en dessous. La règle de cible tient parce qu'elle ajuste à la fois sa cible (vers les signaux lents) et sa vitesse (0.23 de l'écart à $\lambda = 1\,000$) aux coûts.
\end{solution}

\begin{solution}{exo:rs:portfolio-construction-ii:8}
L'écart (0.67 contre 0.57) vaut un quart de l'erreur type d'un ratio de Sharpe sur 7.9 ans (environ 0.4). Il faut changer pour des raisons
mesurées avec précision (stabilité, concentration, rotation), et non sur cette comparaison.
\end{solution}

\begin{solution}{pb:rs:portfolio-construction-ii:1}
\begin{enumerate}
\item Moyenne--variance 0.029~\% et 0.33~\%~; covariance échantillonnale 0.028~\% et 0.24~\%~; robuste 0.028~\% et 0.12~\%~; Black--Litterman 0.031~\% et 0.11~\%.
\item 23.9~\%, 23.4~\%, 22.4~\% et 24.4~\% du capital, dans un sens.
\item Elles ignorent l'alpha~; elles ne peuvent pas non plus l'utiliser.
\item Les plus grands sauts entre titres quasi interchangeables~; et non la réponse moyenne à la prévision.
\item $\mu = [(\tau\Sigma)^{-1} + P^\top\Omega^{-1}P]^{-1}[(\tau\Sigma)^{-1}\Pi + P^\top\Omega^{-1}q]$~; $\tau$ fixe le poids de l'incertitude a priori relativement à celle des vues.
\item Volatilité de 18.4~\% et de 17.8~\%~; la \omterm{def:rs:portfolio-construction-ii:hrp}{parité de risque hiérarchique} (HRP) détient au plus 5.8~\% dans un même titre.
\item Elle ordonne les actifs par regroupement et répartit le capital selon les inverses des variances des sous-blocs, sans jamais inverser la matrice.
\item Les investisseurs réticents à l'effet de levier surpaient les actifs risqués, laissant aux actifs plus sûrs des rendements ajustés du risque plus élevés aux investisseurs qui peuvent recourir au levier (Asness, Frazzini et
Pedersen).
\item 0.57, 0.77, 0.61, 0.72, 0.65~; 0.67, 0.66, 0.65~; erreur type d'environ 0.4.
\item Stabilité, concentration et coûts diffèrent~; les écarts de performance ne sont pas détectables.
\item La dérive persistante ($\phi = 0.00137$ par jour), la dérive post-annonce (1/60), le retournement d'un jour (1).
\item Il poursuit le retournement chaque jour, en renouvelant 2.11 fois le capital par jour pour un coût de 56.6~\% par an.
\item Les signaux lents à poids presque plein, le retournement à 0.56~; elle échange 0.56 de l'écart chaque jour.
\item \textbf{Résultat nommé.} Par point de base~: 0.029~\% (moyenne--variance), 0.028~\% (échantillon et robuste), 0.031~\% (Black--Litterman), 0 pour les règles fondées sur le risque~; à
$\lambda = 100$, le \omterm{def:rs:portfolio-construction-ii:aim}{portefeuille cible} dégage un ratio de Sharpe net de 4.06 contre $-0.51$ pour la réoptimisation quotidienne.
\item Pour $\lambda = 1\,000$, la règle de cible nette 3.15, les autres 2.55 ou moins~; pour $10\,000$, seule la règle de cible est positive (1.61).
\item Elle échange rarement, ce qui est peu coûteux lorsque les coûts sont faibles, mais chaque saut mensuel est grand, et les coûts quadratiques pénalisent les grandes transactions.
\item Le bruit qui se renouvelle chaque jour s'apparente à un signal rapide~: la règle de cible le sous-pondérerait et gagnerait par rapport aux autres.
\item La règle de cible, avec des taux de décroissance estimés signal par signal et un $\lambda$ estimé à partir des exécutions propres à la firme.
\item À partir de l'impact sur les prix des transactions de la firme (chapitre~27)~: régresser le coût sur la taille de la transaction rapportée au volume et à la volatilité.
\item Le chemin~: avec les coûts, la manière dont un portefeuille se déplace détermine davantage sa performance nette que le portefeuille vers lequel il se déplace.
\end{enumerate}
\end{solution}

\begin{solution}{iq:rs:portfolio-construction-ii:1}
L'\omterm{def:rs:portfolio-construction-i:mv}{optimisation moyenne--variance} appliquée aux prévisions brutes donne des portefeuilles extrêmes et instables. Black--Litterman part des rendements qui rendent le marché optimal et s'en écarte seulement autant que
les vues, pondérées par leur confiance, le justifient~: une distribution a posteriori bayésienne dont l'optimum est le marché plus des inclinaisons.
\end{solution}

\begin{solution}{iq:rs:portfolio-construction-ii:2}
Allouer le capital de sorte que chaque classe d'actifs contribue également au risque, en ignorant les rendements espérés (supposés à peu près proportionnels au risque). Pour atteindre des rendements proches de ceux des actions, les
classes à faible risque, surtout obligataires, doivent être employées avec effet de levier~; l'argument en faveur est que les investisseurs réticents à l'effet de levier laissent aux actifs plus sûrs des rendements ajustés du risque plus élevés.
\end{solution}

\begin{solution}{iq:rs:portfolio-construction-ii:3}
Il retranche le pire cas de l'alpha sur un \omterm{def:rs:portfolio-construction-ii:robust}{ensemble d'incertitude}, ce qui pénalise la concentration le long des directions incertaines et amortit les basculements. Dimensionner l'ensemble à partir
de l'erreur estimée de la prévision (un nombre d'erreurs types, ou un quantile de la loi du chi-deux), et vérifier la stabilité et la performance réalisée.
\end{solution}

\begin{solution}{iq:rs:portfolio-construction-ii:4}
Regrouper les actifs selon la distance de corrélation, les ordonner, et répartir le capital de façon récursive selon les inverses des variances des moitiés. L'utiliser lorsque la matrice de covariance est
mal conditionnée ou singulière, et lorsque la stabilité compte davantage que l'utilisation d'un alpha.
\end{solution}

\begin{solution}{iq:rs:portfolio-construction-ii:5}
Échanger partiellement vers un \omterm{def:rs:portfolio-construction-ii:aim}{portefeuille cible} qui sous-pondère le signal rapide d'un facteur $1/(1 + \phi a/\gamma)$ (Gârleanu et Pedersen)~: viser devant la cible, échanger
une partie du chemin.
\end{solution}

\begin{solution}{iq:rs:portfolio-construction-ii:6}
Contributions égales $w_1(\Sigma w)_1 = w_2(\Sigma w)_2$ avec $w_1 + w_2 = 1$~: $w_1^2\sigma_1^2 + w_1w_2\rho\sigma_1\sigma_2 = w_2^2\sigma_2^2 + w_1w_2\rho\sigma_1\sigma_2$, donc
$w_1\sigma_1 = w_2\sigma_2$ quel que soit $\rho$~: poids inverses de la volatilité, $w_1 = \sigma_2/(\sigma_1 + \sigma_2)$.
\end{solution}
